\section{Conclusion}
\label{sec:conclusion}

CoMem opens a new systems axis for repeated-query long-context inference:
\emph{how much transformer depth should a reusable document object prepay?}
Its single-residual interface makes that axis explicit, while the matched
$j{=}0$ endpoint separates depth reuse from retrieval and packing. On
Qwen3-8B, prepaying 12 of 36 layers yields a $1.403\times$ selected-pack Read
speedup at a 3.12-point RULER cost; together with bounded selection, it produces
a $64.9\times$ 128k online-prefill operating point. The accompanying storage,
Write, and break-even measurements show when this trade becomes worthwhile in
a repeated-query workload.

The mechanism analysis turns the quality gap into an actionable design result.
A continuous-prefix oracle recovers all 3.12 points, the context--position
factorization identifies lower-layer document context as the dominant tested
factor on multikey retrieval, and a 32-token overlap recovers 6.0 points without
increasing persistent bytes or per-query Read. The equal-latency study adds an
important boundary: BM25 raw replay remains decisively stronger, whereas the
frozen-BGE comparison is unresolved, making selector quality part of the
systems frontier rather than a hidden constant.

These results establish transformer depth as a measurable and tunable dimension
of reusable context, alongside selection budget, stored representation, and
storage placement. CoMem supplies the interface and matched methodology needed
to study that dimension; the next generation of reusable-context systems can
combine it with stronger selectors, contextual writers, and PIC or learned
modular-cache repair.
